This thesis told a story of theory and evidence: we started from a bound using Hoeffding's inequality and a union bound over the hypothesis class, which did not capture the saturation of the empirical error curve for binary decision trees. Then, we took into account that measuring the complexity of the class of functions that decision trees can learn after they have seen the data is more representative of how the algorithm works in practice. By including this in our analysis via Lemma~\ref{lemma: growth}, we obtained a saturating bound. Finally, we noticed that the generalization performance of decision trees is affected by the level of noise in the data. Thus, using more refined concentration inequalities, we obtained a bound that is a function of the lowest risk that a decision tree can attain, which depends on the noise through the Bayes risk. This, combined with the result on the actual complexity of the hypothesis class, gave a bound that saturates and that is sensitive to the noise. 

The main limitation of this thesis is the lack of an analysis of the difference between the selected model that we consider for theoretical results, i.e.\ the Empirical Risk Minimizer, and the model learned in practice by the DecisionTreeClassifier algorithm. In fact, this is a greedy algorithm that does not always select the best possible model. Among the results we derived, only the lower bound in Theorem~\ref{thm: unionbound} and the results from Theorem~\ref{thm: bound emprisk} onwards are specific to the ERM. Analyzing this issue in more depth might be a direction for future work. 

Besides this, although we obtain a bound that saturates and is noise-sensitive, there is clearly room for improving this study and obtaining functions that capture the qualitative behavior of the generalization gap also under other aspects. In this regard, we believe that a promising direction is to study the qualitative descriptive power of bounds based on a {\it localized\/} version of the Rademacher complexity, proposed by~\cite{Bartlett_2005}. Another potential development is to consider other measures of complexity of the hypothesis class, not necessarily based on the maximum depth. 

Finally, we hope that the framework and approach proposed in this work might be applied to other algorithms in the future. Specifically, a natural next step would be Random Fourier Features, while we propose neural networks trained with Stochastic Gradient Descent as a more ambitious one. 